\documentclass[10pt,a4paper]{article}
\usepackage[margin=23mm]{geometry}
\usepackage[T1]{fontenc}
\usepackage{lmodern}
\usepackage{amsmath,amssymb,booktabs,microtype}
\usepackage[hidelinks]{hyperref}
\setlength{\parindent}{0pt}
\setlength{\parskip}{4pt}
\title{Problem 00000000025:\\Two incompatible assertions about exceptional zeros}
\author{lizaixi01}
\date{24 October 2026}
\begin{document}
\maketitle
\vspace{-15pt}

\textbf{Scope and result.}
The printed conjecture says both that two non-overlapping irreducible characters
of the same Galois extension cannot simultaneously have zeros in the common
window, and that an exceptional zero exists \emph{if and only if the associated
character is real}. We read the latter as the printed assertion for each
associated character. These two assertions are incompatible. The proof uses
two distinct, nonprincipal, primitive real characters of conductor eight.
It asserts no actual existence of a Siegel zero. The usual correct statement
that a possible exceptional character must be real is unaffected, as are the
classical conditional repulsion theorems.

\textbf{Actual field and characters.}
Let $K=\mathbb Q(\zeta_8)$. The native Lean construction uses
\texttt{CyclotomicField} with parameters $8$ and $\mathbb Q$.
It proves that $K/\mathbb Q$ is Galois of degree four, constructs its genuine
automorphism group $G$, and uses the native isomorphism
$G\simeq(\mathbb Z/8\mathbb Z)^\times$ given by the action on $\zeta_8$.
The native number-field discriminant formula gives
$\operatorname{discr}(K)=256$, so $D=|\operatorname{discr}(K)|>1$.
The two actual multiplicative characters are extended by zero on nonunits:
\[
\begin{array}{c|rrrr|c}
u&1&3&5&7&\text{conductor}\\\hline
\chi_A(u)&1&-1&-1&1&8\\
\chi_B(u)&1&1&-1&-1&8
\end{array}
\]
Their multiplicativity is checked on all units. Neither can factor through
modulus four: $5\equiv1\pmod4$, but both take value $-1$ at $5$.
The actual conductor divides eight, and this excludes all its proper divisors.

Transport to $G$ gives the genuine complex one-dimensional representations
$\rho_\chi(g)z=\chi(g)z$. Their native trace characters are exactly these
transported characters. Every subrepresentation is zero or the entire complex
line, so both are irreducible. At the automorphism corresponding to $3$, their
actions are $-1$ and $1$; every intertwiner between them is zero.
The actual normalized character inner product is also computed:
\[
 \frac1{|G|}\sum_{g\in G}\operatorname{tr}(\rho_A(g))
       \overline{\operatorname{tr}(\rho_B(g))}
 =\frac{1-1+1-1}{4}=0.
\]
Thus ``non-overlapping'' is verified by genuine inequivalence, vanishing
intertwiners and orthogonality, and both characters are real-valued.

\textbf{The contradiction.}
For the same arbitrary $c>0$ and this actual $D$, put
\[
 Z_c(\chi)\;:\Longleftrightarrow\;
 \exists\beta\in(1-c/\log D,1),\quad L(\beta,\chi)=0.
\]
This window is nonempty. The printed sufficiency gives both $Z_c(\chi_A)$ and
$Z_c(\chi_B)$, whereas its pairwise exclusion gives their negated conjunction.
The proof works for every function family, with the displayed actual
zero predicate, and then explicitly instantiates Mathlib's analytically
continued \texttt{DirichletCharacter.LFunction}. No zeros or zero sets are
assigned. It refutes only the conjunction of these two necessary printed
clauses; a weaker global assertion about a possible real character is a
different statement.

\newpage
\textbf{Full and depleted Artin factors.}
Milne's introductory cyclotomic identification is stated for $L_S$, with
the ramified primes omitted \cite[pp.~7--9]{milne}. We keep that distinction.
For $K=\mathbb Q(\zeta_8)$ the only ramified rational prime is $2$.
The Lean proof constructs the actual prime ideal $(\zeta_8-1)$ in the native
ring of integers, proves that it lies over $(2)$, computes its ramification
index as four, and proves that its actual inertia subgroup is all of $G$.
Since both representations take value $-1$ at $5$, their genuine
inertia-fixed subspaces are zero. Consequently the full local Artin factor
at $2$ is
\[
 \det(1-2^{-s}\operatorname{Frob}_2\mid V^{I_2})^{-1}=1.
\]
Lean proves this determinant identity for every endomorphism of the actual
inertia-fixed space. On the Dirichlet side $\chi_A(2)=\chi_B(2)=0$, and their
actual local Euler factors are likewise one. Thus adding the ramified factor
does not change either function in this case. The full Artin convention uses
the inertia-fixed determinant at every finite prime \cite[pp.~4--5]{block}.

At odd primes the usual cyclotomic Frobenius is the corresponding unit
class. The standard equality of the resulting Artin/primitive-Dirichlet
Euler products and their analytic continuations is the explicitly disclosed
analytic recognition boundary. A general Artin Euler-product theory and that
general continuation-identification theorem are not reconstructed in Lean.
The formal universal-function-family contradiction already applies to any
such actual Artin family; the native Dirichlet instance additionally supplies
fully specified analytic functions. No analytic recognition theorem is
inserted as an extra Lean premise.

For comparison, the standard zero-free-region theorem leaves \emph{at most}
one possible exceptional real character, conditional on such a zero existing
\cite[Theorem~12.3, p.~119; p.~123]{kouk}. It does not give a zero for every
real character. The present disproof neither decides existence of Siegel
zeros nor contests this standard theorem, its constants, or GRH.

\textbf{Formal coverage and reproduction.}
The complete source is \texttt{lean/Main.lean}, namespace \texttt{TLMC25},
ending in \texttt{complete\_counterexample}. It contains the actual field,
Galois isomorphism, discriminant, characters and native conductors,
representations and traces, all subrepresentations, inner product,
intertwiners, ramification and inertia, local determinants, actual
$L$-function zero conditions, and the quantified contradiction for every
$c>0$. The complete original English and Chinese source is in
\texttt{SOURCE.md}. The project pins Lean~4.33.0 and Mathlib commit
\texttt{db584cd6d46c92f209a44c0f1c829460d327499d}.
From the \texttt{lean} directory, the four independent checks are:
\begin{verbatim}
lean --version
lake build +Main
lake env lean -DwarningAsError=true Check.lean
lake env leanchecker --fresh Main
\end{verbatim}
\texttt{verify.sh} runs the last three commands. The exhaustive check prints
ordinary axioms for every owned environment declaration, including generated
names, and checks exact namespace enumeration coverage. Only
\texttt{propext}, \texttt{Classical.choice} and \texttt{Quot.sound} are allowed.
The source-bound original command logs, job/result data, all actual thread
mask samples, qualification receipts and hash manifest accompany the package.
Production timings are not controlled causal efficiency evidence.

\begin{thebibliography}{9}\small
\bibitem{milne} J.~S.~Milne, \emph{Class Field Theory}, v4.03,
6 August 2020, Introduction, pp.~7--9.
\url{https://www.jmilne.org/math/CourseNotes/CFT.pdf}
\bibitem{block} A.~Block, \emph{Artin L-Functions}, May 2017,
\S3, pp.~4--5.
\url{https://www.columbia.edu/~abb2190/ArtinLFunctions.pdf}
\bibitem{kouk} D.~Koukoulopoulos, \emph{The Distribution of Prime Numbers},
author's preliminary version, Theorem~12.3 and exceptional-character discussion,
pp.~119, 123.
\url{https://dms.umontreal.ca/~koukoulo/documents/publications/primes.pdf}
\end{thebibliography}
\end{document}
